\scanpage{60}
$\operatorname{int}(\sigma_\infty),\operatorname{int}\sigma\tau_\infty)$,
\inserted{Now in this case $g_0=\tau_\infty$ resp. $g_0=\sigma_\infty$, \deleted{thus} \inserted{one} finds $r'=(i,\omega,\varepsilon')$ resp. $r'=(i,\omega',\varepsilon')$, thus}
\[
(39)\qquad\begin{cases}
\psi_i^{\omega,\varepsilon'}=\tau\circ\psi_i^{\omega,\varepsilon}=\psi_i^{\omega,\varepsilon}\circ\tau_\infty\\[5pt]
\psi_i^{\omega',\varepsilon'}=\underbrace{\sigma_i}_{\substack{\text{transposition}\\\in\mathcal G_J\text{ which}\\\text{fixes }i}}\circ\psi_i^{\omega,\varepsilon}=\psi_i^{\omega,\varepsilon}\circ\sigma_\infty
\end{cases}
\]
thus by composing
\[
(40)\qquad\psi_i^{\omega',\varepsilon\deleted{\unsure{'}}}=\deleted{\varepsilon}\underbrace{\widetilde\sigma_i}_{\substack{\tau\sigma_i\\\text{reflection}\\\text{about the}\\\text{great circle}\\\text{passing through}\\P_\varepsilon,P_{\varepsilon'},Q_i}}\circ\psi_i^{\omega,\varepsilon}=\psi_i^{\omega,\varepsilon}\circ\widetilde\sigma_\infty
\]
In all these calculations, probably following old habits, I have favoured the $P_\varepsilon$ and the $Q_i$ as base poin\textsuperscript{ts} -- whereas for the work I have to do, it is the points $R_i^\omega$ that are the most crucial, \unsure{where|for} the calculations made at these poin\textsuperscript{ts} directly capture the phenomena of ``monodromy'' around the ``holes'' $R_i$. I am going to remedy this neglect, by introducing, in
\margin{59}
